% Comandos do formulario 1-dados/dados.tex, inclusao de arquivos e verificacoes. Nao editar.
\makeatletter

% ---------------------------------------------------------
% Mensagens
% ---------------------------------------------------------
\newcount\ppg@nerros
\newcount\ppg@npendencias
\newcommand*\ppg@aviso[1]{\PackageWarningNoLine{Modelo PPGENAP}{#1}}
\newcommand*\ppg@erro[2]{\global\advance\ppg@nerros\@ne\PackageError{Modelo PPGENAP}{#1}{#2}}
% Pendencia: aviso durante o preenchimento; erro na versao para entrega.
\newcommand*\ppg@pendencia[1]{%
  \global\advance\ppg@npendencias\@ne
  \ppg@op@entrega
    {\ppg@erro{#1}{Corrija antes de gerar a versao para entrega.}}%
    {\ppg@aviso{#1}}}

% Versao do modelo e formularios aceitos
\InputIfFileExists{modelo/versao.tex}{}{%
  \ppg@erro{Arquivo modelo/versao.tex nao encontrado}
    {Restaure a pasta modelo/ a partir do ZIP do modelo.}%
  \def\ppg@versaomodelo{?}\def\ppg@datamodelo{?}\def\ppg@formularioatual{1}%
  \def\ppg@formulariominimo{1}}
\PackageInfo{Modelo PPGENAP}{Modelo versao \ppg@versaomodelo\space(\ppg@datamodelo)}

% ---------------------------------------------------------
% Campos do formulario (preenchem os comandos de leitura)
% ---------------------------------------------------------
\newcommand\NomeDoAluno[1]{\def\autor{#1}}
\newcommand\TituloDaDissertacao[1]{\def\titulo{#1}}
\newcommand\NomeDoOrientador[1]{\def\orientador{#1}}
\newcommand\InstituicaoDoOrientador[1]{\def\instorientador{#1}}
\newcommand\NomeDoCoorientador[1]{\def\coorientador{#1}}
\newcommand\Instituicao[1]{\def\instituicao{#1}}
\newcommand\Centro[1]{\def\centro{#1}}
\newcommand\Programa[1]{\def\curso{#1}}
\newcommand\SiglaDoPrograma[1]{\def\siglaprograma{#1}}
\newcommand\TipoDeTrabalho[1]{\def\tipotrabalho{#1}}
\newcommand\NivelDoCurso[1]{\def\nivelcurso{#1}}
\newcommand\NomeDoCurso[1]{\def\nomecurso{#1}}
\newcommand\Grau[1]{\def\grau{#1}}
\newcommand\AreaDoGrau[1]{\def\areagrau{#1}}
\newcommand\CidadeDeDeposito[1]{\def\cidade{#1}}
\newcommand\AnoDeDeposito[1]{\def\ano{#1}}
\newcommand\DataDeAprovacao[1]{\def\datadeaprovacao{#1}}
\newcommand\VersaoDoFormulario[1]{\def\ppg@versaoformulario{#1}}

% Banca: \MembroDaBanca{Nome}{Funcao}{Instituicao}, um por membro.
% O orientador entra automaticamente como primeiro integrante.
\newcount\ppg@nbanca
\newcount\ppg@i
\newcommand\MembroDaBanca[3]{%
  \global\advance\ppg@nbanca\@ne
  \expandafter\def\csname ppg@bn@\the\ppg@nbanca\endcsname{#1}%
  \expandafter\def\csname ppg@bf@\the\ppg@nbanca\endcsname{#2}%
  \expandafter\def\csname ppg@bi@\the\ppg@nbanca\endcsname{#3}}

% Opcoes sim/nao
\newcommand*\ppg@opcao[3]{%
  \edef\ppg@valor{\detokenize{#2}}%
  \edef\ppg@s{\detokenize{sim}}\edef\ppg@n{\detokenize{nao}}\edef\ppg@na{\detokenize{não}}%
  \ifx\ppg@valor\ppg@s
    \expandafter\let\csname ppg@op@#1\endcsname\@firstoftwo
  \else
    \expandafter\let\csname ppg@op@#1\endcsname\@secondoftwo
    \ifx\ppg@valor\ppg@n\else\ifx\ppg@valor\ppg@na\else
      \expandafter\let\csname ppg@op@#1@invalida\endcsname\@ne
      \ppg@erro{Valor "#2" invalido em \string#3 (arquivo 1-dados/dados.tex)}%
        {Escreva sim ou nao. Enquanto isso, o elemento foi omitido.}%
    \fi\fi
  \fi}
\newcommand\IncluirDedicatoria[1]{\ppg@opcao{dedicatoria}{#1}{\IncluirDedicatoria}}
\newcommand\IncluirAgradecimentos[1]{\ppg@opcao{agradecimentos}{#1}{\IncluirAgradecimentos}}
\newcommand\IncluirEpigrafe[1]{\ppg@opcao{epigrafe}{#1}{\IncluirEpigrafe}}
\newcommand\VersaoParaEntrega[1]{\ppg@opcao{entrega}{#1}{\VersaoParaEntrega}}

% Padroes: elementos opcionais incluidos; modo de preenchimento.
\let\ppg@op@dedicatoria\@firstoftwo
\let\ppg@op@agradecimentos\@firstoftwo
\let\ppg@op@epigrafe\@firstoftwo
\let\ppg@op@entrega\@secondoftwo

% ---------------------------------------------------------
% Carregamento do formulario
% ---------------------------------------------------------
% Campo apagado do formulario: comando de leitura vazio (o campo e verificado depois).
\newcommand*\ppg@padrao[1]{\ifdefined#1\else\def#1{}\fi}
% Compatibilidade do formulario com o modelo (regras em modelo/versao.tex).
\ExplSyntaxOn
\cs_new_protected:Npn \ppg@senumero #1
  { \exp_args:Nno \regex_match:nnTF { \A [0-9]+ \Z } {#1} }
\ExplSyntaxOff
\let\ppg@semdados\z@
\newcommand\ppg@verificarversao{%
  \ifdefined\ppg@versaoformulario\else
    \ppg@aviso{Linha \string\VersaoDoFormulario\space ausente em 1-dados/dados.tex.\MessageBreak
      O formulario foi tratado como versao \ppg@formularioatual.\MessageBreak
      Restaure a linha \string\VersaoDoFormulario{\ppg@formularioatual} no inicio do arquivo}%
    \let\ppg@versaoformulario\ppg@formularioatual
  \fi
  \ppg@senumero\ppg@versaoformulario{}{%
    \ppg@erro{Valor "\detokenize\expandafter{\ppg@versaoformulario}" invalido em \string\VersaoDoFormulario\MessageBreak
      (arquivo 1-dados/dados.tex). Nao altere essa linha;\MessageBreak
      o valor do formulario atual e \ppg@formularioatual}%
      {O formulario foi tratado como versao \ppg@formularioatual.}%
    \let\ppg@versaoformulario\ppg@formularioatual}%
  \PackageInfo{Modelo PPGENAP}{Formulario 1-dados/dados.tex: versao \ppg@versaoformulario}%
  \ifnum\ppg@versaoformulario>\ppg@formularioatual\relax
    \ifnum\ppg@formulariominimo=\ppg@formularioatual\relax
      \edef\ppg@aceitos{o formulario versao \ppg@formularioatual}%
    \else
      \edef\ppg@aceitos{formularios \ppg@formulariominimo\space a \ppg@formularioatual}%
    \fi
    \ppg@erro{1-dados/dados.tex usa o formulario versao \ppg@versaoformulario,\MessageBreak
      mais novo que este modelo (versao \ppg@versaomodelo, que aceita\MessageBreak
      \ppg@aceitos). Atualize modelo/, Principal.tex\MessageBreak
      e latexmkrc a partir do ZIP mais recente do modelo}%
      {Seus arquivos pessoais (pastas 1 a 6) nao precisam ser alterados.}%
  \else\ifnum\ppg@versaoformulario<\ppg@formulariominimo\relax
    \ppg@erro{1-dados/dados.tex usa o formulario versao \ppg@versaoformulario,\MessageBreak
      que este modelo (versao \ppg@versaomodelo) nao aceita mais.\MessageBreak
      Copie seus dados para o 1-dados/dados.tex do ZIP atual}%
      {Veja no LEIA-ME como atualizar.}%
  \else\ifnum\ppg@versaoformulario<\ppg@formularioatual\relax
    \ppg@aviso{1-dados/dados.tex usa o formulario versao \ppg@versaoformulario\space(anterior).\MessageBreak
      Ele funciona; itens novos usam o valor padrao.\MessageBreak
      Veja no LEIA-ME como atualizar}%
  \fi\fi\fi}
\newcommand\CarregarDados{%
  \InputIfFileExists{1-dados/dados.tex}{}{%
    \let\ppg@semdados\@ne
    \ppg@erro{Arquivo 1-dados/dados.tex nao encontrado}%
      {Restaure o arquivo de dados a partir do ZIP do modelo.}}%
  \ifnum\ppg@semdados=\z@\ppg@verificarversao\fi
  \ppg@padrao\autor\ppg@padrao\titulo\ppg@padrao\orientador
  \ppg@padrao\instorientador\ppg@padrao\coorientador\ppg@padrao\instituicao
  \ppg@padrao\centro\ppg@padrao\curso\ppg@padrao\siglaprograma
  \ppg@padrao\tipotrabalho\ppg@padrao\nivelcurso\ppg@padrao\nomecurso
  \ppg@padrao\grau\ppg@padrao\areagrau\ppg@padrao\cidade\ppg@padrao\ano
  \ppg@padrao\datadeaprovacao}

% ---------------------------------------------------------
% Verificacoes (executadas no inicio do documento)
% ---------------------------------------------------------
\newcommand*\ppg@sevazio[1]{%
  \if\relax\detokenize\expandafter{#1}\relax
    \expandafter\@firstoftwo\else\expandafter\@secondoftwo\fi}
\newcommand*\ppg@seigual[2]{%
  \edef\ppg@a{\detokenize\expandafter{#1}}\edef\ppg@b{\detokenize{#2}}%
  \ifx\ppg@a\ppg@b\expandafter\@firstoftwo\else\expandafter\@secondoftwo\fi}
% #1 comando de leitura, #2 campo do formulario, #3 texto de exemplo
\newcommand*\ppg@obrigatorio[3]{%
  \ppg@sevazio{#1}%
    {\ppg@erro{Campo \string#2 vazio (arquivo 1-dados/dados.tex)}%
      {Este campo e obrigatorio. Preencha o texto entre as chaves.}}%
    {\ppg@seigual{#1}{#3}%
      {\ppg@pendencia{Campo \string#2 ainda com o texto de exemplo\MessageBreak
        (arquivo 1-dados/dados.tex)}}{}}}
\newcommand*\ppg@exemplomembro{Prof.(a) Dr.(a) Nome do(a) Membro}
\newcommand\ppg@verificarcampos{%
  \ppg@obrigatorio\autor\NomeDoAluno{NOME COMPLETO DO(A) AUTOR(A)}%
  \ppg@obrigatorio\titulo\TituloDaDissertacao{TÍTULO DA DISSERTAÇÃO: SUBTÍTULO, SE HOUVER}%
  \ppg@obrigatorio\orientador\NomeDoOrientador{Prof.(a) Dr.(a) Nome do(a) Orientador(a)}%
  \ppg@obrigatorio\instorientador\InstituicaoDoOrientador{}%
  \ppg@obrigatorio\cidade\CidadeDeDeposito{}%
  \ppg@obrigatorio\ano\AnoDeDeposito{}%
  \ppg@obrigatorio\instituicao\Instituicao{}%
  \ppg@obrigatorio\centro\Centro{}%
  \ppg@obrigatorio\curso\Programa{}%
  \ppg@obrigatorio\tipotrabalho\TipoDeTrabalho{}%
  \ppg@obrigatorio\nivelcurso\NivelDoCurso{}%
  \ppg@obrigatorio\nomecurso\NomeDoCurso{}%
  \ppg@obrigatorio\grau\Grau{}%
  \ppg@obrigatorio\areagrau\AreaDoGrau{}%
  \ppg@seigual\coorientador{Prof.(a) Dr.(a) Nome do(a) Coorientador(a)}%
    {\ppg@pendencia{Campo \string\NomeDoCoorientador\space ainda com o texto de exemplo.\MessageBreak
      Sem coorientacao, deixe as chaves vazias: \string\NomeDoCoorientador{}}}{}%
  \ppg@sevazio\datadeaprovacao
    {\ppg@pendencia{Data de aprovacao vazia (\string\DataDeAprovacao, arquivo 1-dados/dados.tex)}}%
    {\ppg@seigual\datadeaprovacao{dia de mês de ano}%
      {\ppg@pendencia{Data de aprovacao ainda com o texto de exemplo\MessageBreak
        (\string\DataDeAprovacao, arquivo 1-dados/dados.tex)}}{}}%
  \ifnum\ppg@nbanca=\z@
    \ppg@pendencia{A banca tem apenas o orientador. Acrescente os demais\MessageBreak
      membros com \string\MembroDaBanca\space em 1-dados/dados.tex}%
  \fi
  \ppg@i\z@
  \@whilenum\ppg@i<\ppg@nbanca\do{%
    \advance\ppg@i\@ne
    \ppg@verificarmembro{\the\ppg@i}}%
  \IfFileExists{1-dados/ficha_catalografica.pdf}{}{\IfFileExists{ficha_catalografica.pdf}%
    {\ppg@pendencia{ficha_catalografica.pdf esta fora da pasta 1-dados e nao sera usada.\MessageBreak
      Mova o arquivo para a pasta 1-dados}}%
    {\ppg@pendencia{Ficha catalografica ainda nao inserida.\MessageBreak
      Coloque ficha_catalografica.pdf na pasta 1-dados}}}%
  \ppg@verificarlicenca}
% Sem 1-dados/dados.tex, so o erro do arquivo e mostrado (sem um erro por campo).
\newcommand\ppg@verificar{%
  \ifnum\ppg@semdados=\z@\ppg@verificarcampos\fi
  \IfFileExists{4-referencias/referencias.bib}{}{%
    \ppg@erro{Arquivo 4-referencias/referencias.bib nao encontrado}%
      {Sem ele, a lista de referencias e as citacoes nao sao geradas. Restaure o arquivo.}}}
\newcommand*\ppg@verificarmembro[1]{%
  \expandafter\let\expandafter\ppg@mn\csname ppg@bn@#1\endcsname
  \expandafter\let\expandafter\ppg@mf\csname ppg@bf@#1\endcsname
  \expandafter\let\expandafter\ppg@mi\csname ppg@bi@#1\endcsname
  \let\ppg@incompleto\relax
  \ppg@sevazio\ppg@mi{\def\ppg@incompleto{instituicao}}{}%
  \ppg@sevazio\ppg@mf{\def\ppg@incompleto{funcao}}{}%
  \ppg@sevazio\ppg@mn{\def\ppg@incompleto{nome}}{}%
  \ifx\ppg@incompleto\relax
    \edef\ppg@a{\detokenize\expandafter{\ppg@mn}}%
    \edef\ppg@b{\detokenize\expandafter{\orientador}}%
    \ifx\ppg@a\ppg@b
      \ppg@aviso{O orientador aparece tambem em \string\MembroDaBanca.\MessageBreak
        Ele ja entra automaticamente; remova a linha repetida}%
    \fi
    \ppg@seigual\ppg@mi{Nome da Instituição por Extenso}%
      {\ppg@pendencia{Membro #1 da banca: instituicao ainda com o texto de exemplo}}{}%
    \edef\ppg@a{\detokenize\expandafter{\ppg@mn}}%
    \edef\ppg@b{\detokenize\expandafter{\ppg@exemplomembro}}%
    \@expandtwoargs\in@{\ppg@b}{\ppg@a}%
    \ifin@\ppg@pendencia{Membro #1 da banca: nome ainda com o texto de exemplo}\fi
  \else
    \ppg@erro{Membro #1 da banca sem \ppg@incompleto\space(\string\MembroDaBanca, arquivo 1-dados/dados.tex)}%
      {Use \string\MembroDaBanca{Nome}{Funcao}{Instituicao}, com os tres campos preenchidos.}%
  \fi}
\AtBeginDocument{\ppg@verificar}

% ---------------------------------------------------------
% Folha de aprovacao: orientador e membros cadastrados
% ---------------------------------------------------------
\newcommand\ImprimirBanca{%
  \vfill
  \begin{center}\orientador\par Orientador(a)\par\instorientador\end{center}%
  \ppg@i\z@
  \@whilenum\ppg@i<\ppg@nbanca\do{%
    \advance\ppg@i\@ne
    \vfill
    \begin{center}%
      \csname ppg@bn@\the\ppg@i\endcsname\par
      \csname ppg@bf@\the\ppg@i\endcsname\par
      \csname ppg@bi@\the\ppg@i\endcsname
    \end{center}}}

% ---------------------------------------------------------
% Licenca Creative Commons (opcional, parte inferior da folha de rosto)
% Textos e endereco juridico conforme orientacoes do SISBI/UFRN.
% Nenhuma modalidade e escolhida automaticamente.
% ---------------------------------------------------------
\IfFileExists{ccicons.sty}{\RequirePackage{ccicons}}{}
\newcommand\LicencaCreativeCommons[1]{\ppg@opcao{licenca}{#1}{\LicencaCreativeCommons}}
\newcommand\ModalidadeDaLicenca[1]{\def\ppg@licmodalidade{#1}}
\let\ppg@op@licenca\@secondoftwo
\def\ppg@licmodalidade{}
\let\ImprimirLicenca\@empty
\newcommand*\ppg@licopcoes{CC BY, CC BY-SA, CC BY-ND, CC BY-NC, CC BY-NC-SA ou CC BY-NC-ND}
% \ppg@deflic{CODIGO}{caminho na URL}{icones ccicons}{nome}{descricao}
% Aceita tambem as grafias sem "CC", com "CC-" e com "4.0" no final.
\newcommand\ppg@deflic[5]{%
  \@namedef{ppg@lic@#1}{\def\ppg@liccodigo{CC #1 4.0}\def\ppg@licurl{#2}%
    \def\ppg@licicones{cc#3}\def\ppg@licnome{#4}\def\ppg@lictexto{#5}}%
  \@for\ppg@p:={CC,CC-}\do{\ppg@liccopia{\ppg@p#1}{#1}\ppg@liccopia{\ppg@p#14.0}{#1}}%
  \ppg@liccopia{#14.0}{#1}}
\newcommand*\ppg@liccopia[2]{%
  \expandafter\let\csname ppg@lic@#1\expandafter\endcsname\csname ppg@lic@#2\endcsname}
\ppg@deflic{BY}{by}{by}{Atribuição 4.0 Internacional}{Permite que outros
  distribuam, remixem, adaptem e desenvolvam seu trabalho, mesmo comercialmente,
  desde que creditem a você pela criação original.}
\ppg@deflic{BY-SA}{by-sa}{bysa}{Atribuição-CompartilhaIgual 4.0 Internacional}{Permite
  que outros remixem, adaptem e desenvolvam seu trabalho, mesmo para fins
  comerciais, desde que creditem a você e licenciem suas novas criações sob os
  mesmos termos.}
\ppg@deflic{BY-ND}{by-nd}{bynd}{Atribuição-SemDerivações 4.0 Internacional}{Permite
  que outros reutilizem o trabalho para qualquer finalidade, inclusive
  comercialmente, no entanto, ele não pode ser compartilhado com outros de forma
  adaptada, e o crédito deve ser fornecido a você.}
\ppg@deflic{BY-NC}{by-nc}{bync}{Atribuição-NãoComercial 4.0 Internacional}{Permite
  que outros remixem, adaptem e construam seu trabalho de forma não comercial e,
  embora seus novos trabalhos também devam reconhecê-lo e não serem comerciais,
  eles não precisam licenciar seus trabalhos derivados nos mesmos termos.}
\ppg@deflic{BY-NC-SA}{by-nc-sa}{byncsa}{Atribuição-NãoComercial-CompartilhaIgual 4.0
  Internacional}{Permite que outros remixem, adaptem e desenvolvam seu trabalho de
  forma não comercial, desde que creditem você e licenciem suas novas criações
  sob os mesmos termos.}
\ppg@deflic{BY-NC-ND}{by-nc-nd}{byncnd}{Atribuição-NãoComercial-SemDerivações 4.0
  Internacional}{Esta é a mais restritiva das seis licenças principais, só
  permitindo que outros façam download dos seus trabalhos e os compartilhem desde
  que atribuam crédito a você, mas sem que possam alterá-los de nenhuma forma ou
  utilizá-los para fins comerciais.}
\newcommand\ppg@verificarlicenca{%
  \edef\ppg@m{\expandafter\zap@space\detokenize\expandafter{\ppg@licmodalidade} \@empty}%
  \uppercase\expandafter{\expandafter\def\expandafter\ppg@m\expandafter{\ppg@m}}%
  \ppg@op@licenca
    {\ppg@sevazio\ppg@m
      {\global\let\ImprimirLicenca\ppg@licfalta
       \ppg@erro{Licenca Creative Commons ativada sem modalidade.\MessageBreak
         Escreva em \string\ModalidadeDaLicenca\space(arquivo 1-dados/dados.tex)\MessageBreak
         a modalidade escolhida por voce: \ppg@licopcoes.\MessageBreak
         Sem licenca, use \string\LicencaCreativeCommons{nao}}%
         {O modelo nao escolhe a modalidade por voce. Enquanto isso, a licenca nao e exibida.}}%
      {\ifcsname ppg@lic@\ppg@m\endcsname
         \csname ppg@lic@\ppg@m\endcsname
         \global\let\ImprimirLicenca\ppg@licbloco
         \PackageInfo{Modelo PPGENAP}{Licenca \ppg@liccodigo\space na folha de rosto}%
         \ifcsname\ppg@licicones\endcsname\else
           \ppg@aviso{Pacote ccicons nao encontrado: a licenca sera\MessageBreak
             exibida sem os simbolos Creative Commons}%
         \fi
       \else
         \global\let\ImprimirLicenca\ppg@licfalta
         \ppg@erro{Modalidade de licenca "\detokenize\expandafter{\ppg@licmodalidade}" invalida\MessageBreak
           (\string\ModalidadeDaLicenca, arquivo 1-dados/dados.tex).\MessageBreak
           Use uma destas: \ppg@licopcoes}%
           {Copie a modalidade exatamente como na lista. Enquanto isso, a licenca nao e exibida.}%
       \fi}}%
    {\ifdefined\ppg@op@licenca@invalida\else\ppg@sevazio\ppg@m{}{\ppg@aviso{\string\ModalidadeDaLicenca\space preenchida, mas
       \string\LicencaCreativeCommons{nao}:\MessageBreak a licenca nao sera exibida}}\fi}}
% Bloco impresso: simbolos, codigo, texto (espaco simples) e endereco juridico.
\newcommand\ppg@licbloco{%
  \begingroup
  \setlength{\parindent}{0pt}%
  \fontsize{11}{13.2}\selectfont
  \smallskip
  \ifcsname\ppg@licicones\endcsname
    {\fontsize{20}{20}\selectfont\csname\ppg@licicones\endcsname}\quad
  \fi
  \textbf{\ppg@liccodigo}\par
  \smallskip
  Esta obra está licenciada com uma licença Creative Commons \ppg@licnome.
  \ppg@lictexto\par
  \edef\ppg@a{\noexpand\url{https://creativecommons.org/licenses/\ppg@licurl/4.0/legalcode}}%
  Link dessa licença: \ppg@a\par
  \endgroup}
\newcommand\ppg@licfalta{%
  \begingroup
  \setlength{\parindent}{0pt}%
  \fontsize{11}{13.2}\selectfont
  \textbf{[FALTA: modalidade da licença Creative Commons em 1-dados/dados.tex]}\par
  \endgroup}

% ---------------------------------------------------------
% Inclusao de arquivos do aluno
% ---------------------------------------------------------
\newcommand*\ppg@falta[2]{%
  \clearpage\thispagestyle{empty}\noindent\textbf{[FALTA: #1]}\par\medskip
  \noindent Arquivo #2 nao encontrado.\clearpage}
% \IncluirOpcional{opcao}{arquivo}{Nome}: incluido se a opcao for sim.
\newcommand\IncluirOpcional[3]{%
  \csname ppg@op@#1\endcsname
    {\InputIfFileExists{#2.tex}{}{\ppg@pendencia{#3 marcado como sim em 1-dados/dados.tex,\MessageBreak
      mas o arquivo #2.tex nao foi encontrado; elemento omitido}}}%
    {}}
% \IncluirObrigatorio{arquivo}{Nome}: erro e marca no PDF se faltar.
\newcommand\IncluirObrigatorio[2]{%
  \InputIfFileExists{#1.tex}{}{%
    \ppg@erro{#2: arquivo #1.tex nao encontrado}{Restaure o arquivo ou confira o nome.}%
    \ppg@falta{#2}{#1.tex}}}
% \IncluirSeExistir{arquivo}{Nome}: omitido, com aviso, se faltar.
\newcommand\IncluirSeExistir[2]{%
  \InputIfFileExists{#1.tex}{}{\ppg@aviso{#2: arquivo #1.tex nao encontrado; elemento omitido}}}
% Secoes: listadas em 3-secoes/ordem-das-secoes.tex.
\newcount\ppg@nsecoes
\newcommand\IncluirSecao[1]{%
  \global\advance\ppg@nsecoes\@ne
  \InputIfFileExists{3-secoes/#1.tex}{}{%
    \ppg@erro{Secao "#1" listada em 3-secoes/ordem-das-secoes.tex,\MessageBreak
      mas o arquivo 3-secoes/#1.tex nao existe}%
      {Confira o nome do arquivo (sem .tex) ou apague a linha \string\IncluirSecao{#1}.}%
    \ppg@falta{Secao #1}{3-secoes/#1.tex}}}
\newcommand\IncluirSecoes{%
  \IfFileExists{3-secoes/ordem-das-secoes.tex}{%
    % ORDEM DAS SEÇÕES — cada linha inclui um arquivo desta pasta, nesta ordem.
% Nova seção: crie o arquivo nesta pasta (ex.: 6-produto-educacional.tex)
%   e acrescente a linha \IncluirSecao{6-produto-educacional}
% Remover: apague a linha. Reorganizar: mude a ordem das linhas.
\IncluirSecao{1-introducao}
\IncluirSecao{2-fundamentacao_teorica}
\IncluirSecao{3-metodologia}
\IncluirSecao{4-resultados_discussao}
\IncluirSecao{5-consideracoes_finais}
%
    \ifnum\ppg@nsecoes=\z@
      \ppg@erro{Nenhuma secao listada em 3-secoes/ordem-das-secoes.tex}%
        {Liste as secoes, uma por linha, com \string\IncluirSecao{nome-do-arquivo}.}%
    \fi}{%
    \ppg@erro{Arquivo 3-secoes/ordem-das-secoes.tex nao encontrado}%
      {Restaure o arquivo: ele define quais secoes entram e em que ordem.}%
    \ppg@falta{Ordem das secoes}{3-secoes/ordem-das-secoes.tex}}}

% ---------------------------------------------------------
% Resumo ao final da compilacao
% Na versao para entrega, qualquer problema gera um erro final:
% o PDF nao deve ser depositado. A verificacao automatica nao
% substitui a revisao do(a) orientador(a) e da biblioteca.
% ---------------------------------------------------------
\newcommand\ppg@resumo{%
  % Entrega: referencias ou citacoes indefinidas (?? no PDF). Ignorado enquanto a
  % compilacao anterior ainda nao tinha as referencias do BibTeX (.bbl).
  \ppg@op@entrega{\ifnum\ifdefined\ppg@auxanterior\ppg@auxanterior\else\z@\fi>\z@
    \ifx\@refundefined\relax\else
      \ppg@erro{Ha referencias cruzadas ou citacoes indefinidas (?? no PDF).\MessageBreak
        Compile novamente; se continuar, confira os \string\ref\space e \string\cite}%
        {Uma nova compilacao costuma resolver; se nao, o rotulo ou a chave nao existe.}%
    \fi
  \fi}{}%
  \IfFileExists{\jobname.bbl}%
    {\immediate\write\@auxout{\gdef\string\ppg@auxanterior{1}}}%
    {\immediate\write\@auxout{\gdef\string\ppg@auxanterior{0}}}%
  \ppg@op@entrega
    {\ifnum\ppg@nerros>\z@
       \PackageError{Modelo PPGENAP}{VERSAO PARA ENTREGA COM \the\ppg@nerros\space PROBLEMA(S).\MessageBreak
         Este PDF nao deve ser depositado. Corrija os erros\MessageBreak
         indicados acima e compile novamente}%
         {Para continuar escrevendo, use \string\VersaoParaEntrega{nao} em 1-dados/dados.tex.}%
     \else
       \ppg@aviso{Versao para entrega sem problemas na verificacao automatica.\MessageBreak
         Ela nao garante conformidade integral com a ABNT e a UFRN:\MessageBreak
         revise com o(a) orientador(a) e a biblioteca antes do deposito}%
     \fi}
    {\ppg@aviso{Modo de preenchimento (modelo \ppg@versaomodelo): \the\ppg@nerros\space erro(s)\MessageBreak
       e \the\ppg@npendencias\space pendencia(s). Para gerar o PDF do deposito,\MessageBreak
       use \string\VersaoParaEntrega{sim} em 1-dados/dados.tex}}}
\AtEndDocument{\ppg@resumo}

\makeatother
